% !TeX root = ../教材题库_全部条目_转写工作稿.tex
% Source: XJ-B1-C02-SU-P060-Y1427
% 第2章 / 小结与复习 小结与复习 / 教材 p.60 / PDF 第67页
% OCR未校对：公式与文字均可能有误；未完成数学TeX转写；禁止直接发给学生
% 原题及相邻原解截图（核对用）：../题目原图/XJ-B1-C02-SU-P060-Y1427-01.webp
% 本文件只在转写完成后作为题干使用；OCR 条目尚不是可靠题干。
\begin{bankitem}{第2章 小结与复习 第5题}{第2章 / 小结与复习 小结与复习 / 教材 p.60 / PDF 第67页}{XJ-B1-C02-SU-P060-Y1427}{OCR未校对：公式与文字均可能有误；未完成数学TeX转写；禁止直接发给学生}
\begin{ocrdraft}
[教材p.60 / PDF第67页]
\par
5．不等式可以刻画现实世界中的数量关系．尝试从实际问题中抽象出一元二
\par
次不等式模型，并进行求解，进而解决简单的实际问题．
\end{ocrdraft}
\end{bankitem}
